%% =====================================================================
%%  CHAPTER 4: LITERATURE REVIEW AND FLAGSHIP PAPER   (budget 8--9 pages)
%%  chapters/ch4-literature.tex
%% =====================================================================
\pcChapter{lit}{9}{סקירת ספרות ומאמר הדגל}
\begin{pcguide}[כללי הסמינריון]
רק מאמרים שפיטים (לא בלוגים ולא פודקאסטים); מ-\latin{2023} ואילך (ספרי יסוד מותרים ישנים יותר);
\textbf{מאמר דגל אחד} --- הרלוונטי ביותר, והמקרה הגרוע ביותר: אם דבר אחר לא יעבוד, מיישמים אותו.
כל מקור נפתח ונקרא, עם קישור שעובד. הסקירה מחולקת ל-\rng{3}{4} נושאי ליבה ולא לפי סדר כרונולוגי.
תקציב: \latin{2.5} + \latin{1} + \latin{1.5} + \latin{1} + \latin{2} + \latin{1} עמודים לפי הסעיפים.
\end{pcguide}

\section{סקירה תמטית וביקורתית}\label{sec:review}
הספרות מאורגנת בארבע תמות לפי שרשרת הערך של הפתרון: \textbf{שימור ידע סמוי} (מה צריך לשמר), \textbf{אחזור מבוסס-מקורות
ומיתון הזיות} (איך מעגנים אותו במסמכים, בגרסאות ובסביבה מאובטחת), \textbf{הערכה פדגוגית אוטומטית} (איך הופכים אותו
ליחידות לימוד ומבחנים), ו\textbf{הכשרה טכנית \mbox{תלת-ממדית} ובמציאות מורחבת} (איך מתרגלים אותו). הסינתזה נשענת על
\latin{20} מאמרים שפיטים מהשנים \rng{2023}{2026}, ועל מקור יסוד אחד מ-\latin{2020}. כל מחקר נבחן לפי ארבע שאלות: באיזה הקשר
נבדק, באיזו שיטה, מה בדיוק נמדד, ומה מגבלתו ביחס לסביבה של \latin{IAI}: מסמכים גרסאיים, ידע סמוי, עיבוד מקומי ותרגול
\latin{3D}. ממצאי היצרנים והדוחות המסחריים אינם משמשים כאן, והם מופיעים רק בפרק~\ref{ch:market}.

\subsection{תמה 1: שימור ידע סמוי והעברתו בין דורות}
הסקירה השיטתית של \latin{Igoa-Iraola} ו-\latin{Díez} בוחנת נהלים להעברת ידע ארגוני בעת חילופי דורות, ומסיקה שאין שיטה אחת
מספקת: העברה אפקטיבית משלבת תיעוד, חפיפה אישית, למידה חברתית וחניכה \cite{igoa2024transfer}. מכאן שמאגר מסמכים, גם אם מלא,
אינו תחליף לתהליך שבו הידע מוסבר ומתורגל. \latin{Finkel} ו-\latin{Wurster} חוקרים ידע סמוי בייצור תעשייתי: הם ממפים חסמים
ופרקטיקות קיימות, ומציעים כיווני יישום לכלי ניהול ידע מבוססי \latin{LLM} \cite{finkel2026unlocking}. \latin{Zuin} ועמיתיו
מציגים גישה לגילוי ידע סמוי בהקשר ארגוני באמצעות מודלי שפה \cite{zuin2025tacit}. שלושת המחקרים מתכנסים לעמדה אחת:
ידע סמוי אינו נמצא במסמכים, וכדי להפוך אותו לנגיש יש לחלץ אותו בשיחה מובנית מאנשים.

הביקורת על קבוצה זו כפולה. ראשית, הראיות הן בעיקר מיפוי, סקירה והדגמה: אף אחד מהמחקרים לא מגדיר מדד שבו אפשר לקבוע אם
הטענה שחולצה מן המומחה נכונה, כלומר \textbf{דיוק חילוץ}. שנית, חילוץ ידע בעזרת מודל שפה מציב סיכון הפוך להזיה: מודל המסכם
ריאיון עלול לנסח "ידע" שהמומחה לא אמר. מסקנה לתכן: כל טענה שחולצה תיבדק מול מקור תקף או תסומן כטענה הממתינה לאישור מומחה,
ודיוק החילוץ יהפוך למדד מוצהר (\req{16}, \hyp{1}). ללא מדד כזה אי אפשר להצהיר שהידע נשמר, אלא רק שהוא נאסף.

\subsection{תמה 2: \latin{Grounded RAG} ומיתון הזיות}
הבסיס הרעיוני הוא הצמדת מודל שפה למאגר מסמכים שממנו מאוחזרות ראיות \cite{lewis2020retrieval}; זהו מקור היסוד היחיד שקדם
ל-\latin{2023}. הצורך בעיגון נובע מהזיות: \latin{Ji} ועמיתיו מסכמים את סוגי ההזיה במודלי יצירת טקסט ואת מגבלות שיטות המדידה
שלהן \cite{ji2023survey}, ו-\latin{Barnett} ועמיתיו מתעדים שבע נקודות כשל בהנדסת מערכות \latin{RAG} במקרי בוחן מעשיים,
החל במידע שאינו במאגר וכלה בתשובה שאינה נשלפת מההקשר \cite{barnett2024seven}. מכאן שאי אפשר לתכנן \latin{RAG} כ"קופסה שחורה":
כל נקודת כשל מחייבת בקרה.

ההערכה היא התחום הבוגר ביותר בתמה. \latin{RAGAs} מציע מדדים ללא תשובת אמת לכל שאילתה, ובהם נאמנות לראיות ורלוונטיות
\cite{es2024ragas}; \latin{ARES} משלב שופט אוטומטי מכויל בשיפוט אנושי כדי לבדוק את הרכיבים בנפרד \cite{saad2024ares}.
הביקורת: שופטים אלו נבדקו בעיקר על קורפוסים פתוחים, ולכן אין ודאות שהם אמינים על מסמכי נוהל טכניים בעברית ובאנגלית; בפיילוט
נשתמש בהם רק יחד עם שיפוט אנושי של שני מעריכים. מבחינת מבנה המסמך, \latin{TreeRAG} מציע אחסון היררכי לשיפור האחזור במסמכים ארוכים \cite{tao2025treerag}, ו-\latin{Lukens} ועמיתיו מציעים מסגרת
הערכה לאבחון תקלות המסתמך על מדריכים טכניים \cite{lukens2025manuals}; שניהם מדגישים שמבנה המסמך משפיע על איכות האחזור.

שני מחקרים נוגעים ישירות באילוצי \latin{IAI}. \latin{KnowShiftQA} בוחן את חוסנן של מערכות \latin{RAG} כאשר הידע בספר הלימוד שונה מן הידע הכללי של המודל
\cite{zheng2025knowshift}; זוהי בעיה קרובה לערבוב בין גרסאות, הסיכון המרכזי בתחזוקת נהלים. ומחקרם של
\latin{Zeng} ועמיתיו מראה שמאגר אחזור פרטי עלול לדלוף באמצעות שאילתות מתוכננות \cite{zeng2024good}, ולכן פריסה מקומית
אינה פוטרת מסינון לפי הרשאה ומיומן שאילתות. הפער בקבוצה זו הוא שהמחקרים בוחנים כשלים בנפרד: איש אינו מחבר בין
\textbf{גרסת מסמך}, \textbf{הרשאה} ו\textbf{הימנעות ממענה} במערכת אחת שמופעלת גם על ידע שחולץ מבני אדם.

\subsection{תמה 3: הערכה פדגוגית אוטומטית והפקת שאלות}
\latin{Guinet} ועמיתיו מציעים להעריך מערכת \latin{RAG} באמצעות מבחן רב-ברירה שנוצר אוטומטית מקורפוס המשימה, ומשתמשים בתורת
התגובה לפריט \latin{(IRT)} כדי לזהות שאלות בעלות מידע נמוך ולשפר את המבחן באופן איטרטיבי \cite{guinet2024exam}. \latin{Pradeesh}
ועמיתיו מיישמים \latin{RAG} להפקת שאלות רב-ברירה ותשובותיהן \cite{pradeesh2025mcq}, ו-\latin{Ramkumar} ו-\latin{Rajendran}
מתארים גישה ליצירת תמיכה פדגוגית מעוגנת בחומרי הלימוד \cite{ramkumar2026pedagogical}. כולם מראים שאפשר להפיק שאלות
מקורפוס בעלות נמוכה, אך מדידת איכותן נשענת בעיקר על מדדים טכניים ולא על מחוון פדגוגי עם אישור מומחה. סכנה נוספת עולה
ממחקרם של \latin{Lee} ועמיתיו על עובדי ידע: עובדים המשתמשים בבינה מלאכותית יוצרת דיווחו על ירידה במאמץ החשיבה הביקורתית,
במיוחד כשהם סומכים על הכלי \cite{lee2025impact}. לפיכך חומר הדרכה שנוצר אוטומטית עלול לעודד קבלה לא ביקורתית, ולכן התרגול
יכלול שאלות הנמקה ולא רק זיהוי.

מאמר הדגל, \latin{Hua} ועמיתיו, שייך לתמה זו, ומנותח לעומק בסעיף~\ref{sec:flagship} \cite{hua2025civask}: הוא מציע הערכה יזומה של
עוזר למידה מבוסס \latin{RAG} במקום הסתמכות על תלונות משתמשים. הפער בתמה: אין מחקר שמודד את \textbf{זמן ההפקה} של חומר
מאושר לעומת הכנה ידנית, ואין מחקר שמחייב אישור מומחה כשער לפני פרסום \req{06}.

\subsection{תמה 4: הכשרה טכנית \mbox{תלת-ממדית} ובמציאות מורחבת}
המטא-אנליזה של \latin{Bödding} ועמיתיו, שכללה \latin{53} מחקרים על מציאות מעורבת בהכשרה מקצועית, מצאה השפעה חיובית
על תוצאות התנהגותיות \latin{($d=0.40$)}, קוגניטיביות \latin{($d=0.84$)} ורגשיות \latin{($d=0.65$)}, ואפקטיביות שהייתה ברובה
בלתי תלויה במאפייני הטכנולוגיה עצמה \cite{bodding2025mr}. \latin{Somerville} ועמיתיו סקרו \latin{67} מחקרים על מציאות
מורחבת באימון טייסים בסימולטורים ומצאו אפקט חיובי מובהק, בינוני עד חזק \latin{(0.884)}, במטא-אנליזה של חמישה מחקרים בלבד; הם מציינים
שאף מחקר אינו ניסוי אקראי מבוקר, ושנדרש מחקר נוסף על בחילת סימולטור והעברת ההכשרה לשטח \cite{somerville2025xr}.
\latin{Jalilvand} ועמיתיו משווים מערכות \latin{VR/MR} לאימון בתהליך ייצור בשיטת קבלת החלטות רב-קריטריונית
\cite{jalilvand2024xr}, ו-\latin{Yang} ועמיתיו סוקרים שילוב מציאות מורחבת בתכנון פירוק רכיבי מטוס בסוף חייהם, וממפים חסמי
אימוץ \cite{yang2023xr}.

הביקורת: ראשית, ההשפעה תלויה בתוכן ובשיטת ההוראה יותר מבטכנולוגיה, ולכן ערך ה-\latin{3D} אינו מובן מאליו וצריך
להיבחן מול תרגיל \mbox{דו-ממדי} זהה בתוכן. שנית, בסיס הראיות במחקרים בתחום התעופה צר ואינו מבוסס על ניסויים אקראיים. שלישית ---
וזהו הפער הנדסי --- איש אינו מדווח על עלות ההכנה של הסצנות: כל המחקרים מניחים שהתרחיש כבר קיים, ואינם בוחנים כיצד
יוצרים אותו ומעדכנים אותו כשהנוהל משתנה. כאן נכנס רעיון “הכלי שמייצר את הכלי” של הפרויקט (סעיף~\ref{sec:solconcept}).

טבלה~\ref{tab:themes} מסכמת את הסינתזה: לכל תמה מה מוסכם, מה חסר, ואיזו דרישה נגזרת.
\begin{table}[!htbp]
\centering
\caption{סינתזה ביקורתית לפי תמות: מה ידוע, מה חסר ומה נגזר לפרויקט}\label{tab:themes}
\begin{pctable}{H{0.19},H{0.26},H{0.30},H{0.25}}
  \pchead{תמה & מה מוסכם בספרות & מה חסר או שנוי במחלוקת & נגזרת לתכן}
  \pcrow{ידע סמוי והעברת ידע & ידע סמוי דורש חילוץ מובנה מאנשים ולא רק תיעוד & אין מדד דיוק לחילוץ; סיכון שמודל ינסח מה שהמומחה לא אמר & \req{16}, \hyp{1}: טענה מאושרת רק מול מקור או מומחה}
  \pcrow{\latin{Grounded RAG} והזיות & עיגון במקור מקטין הזיות; קיימים מדדי נאמנות ושופטים אוטומטיים & גרסאות משתנות, הרשאות ודליפה נבחנים בנפרד; שופטים לא נבדקו על נהלים טכניים & \req{02}, \req{03}, \req{05}, \req{09}: עיגון בגרסה, הימנעות ושיפוט אנושי}
  \pcrow{הערכה פדגוגית אוטומטית & אפשר להפיק שאלות מקורפוס בעלות נמוכה & איכות נמדדת טכנית; אין אישור מומחה וזמן הפקה מול ידני; סיכון לקבלה לא ביקורתית & \req{04}, \req{06}, \req{12}: מחוון ואישור מומחה, מדידת זמן}
  \pcrow{הכשרה ב-\latin{3D}/\latin{XR} & השפעה חיובית בינונית עד גדולה, תלוית תוכן והוראה & ראיות צרות בתעופה; עלות ויצירת הסצנות ועדכונן אינם נבחנים & \req{13}, \req{15}, \hyp{4}: סצנה פרמטרית מ-\latin{JSON}; השוואה מול \latin{2D}}
\end{pctable}
\end{table}

\subsection{סינתזה חוצת תמות: מפת הספרות ושיטות ההערכה שאומצו}\label{sec:crossmap}
איור~\ref{fig:litmap} ממקם את עשרים המאמרים בארבע התמות ומסמן בין כל שתי תמות עוקבות את החיבור שלא נבדק במחקרים שנסקרו. הפער אינו בתמה
בודדת אלא בצירופים: כל תמה נחקרת בנפרד, ואיש אינו מחבר אותן לשרשרת אחת מידע-ליחידה-לתרגול.
\begin{figure}[!htbp]
\centering
\pcFig{litmap}{%
\pcAutoFit{\begin{pctikz}[x=1cm,y=1cm]
  \tikzset{th/.style={pcnodeD,text width=2.6cm,minimum height=1.3cm,font=\small\sffamily},
           gp/.style={pcnode,dashed,draw=pcBad,fill=white,text width=3.2cm,font=\scriptsize\sffamily,anchor=north}}
  \node[th] (a) at (0,0) {\tn{תמה 1: ידע סמוי\\[-1pt]\latin{3} מאמרים}};
  \node[th] (b) at (-3.6,0) {\tn{תמה 2: \latin{Grounded RAG}\\[-1pt]\latin{8} מאמרים}};
  \node[th] (c) at (-7.2,0) {\tn{תמה 3: הערכה פדגוגית\\[-1pt]\latin{5} מאמרים}};
  \node[th] (d) at (-10.8,0) {\tn{תמה 4: \latin{3D/XR}\\[-1pt]\latin{4} מאמרים}};
  \foreach \p/\q in {a/b,b/c,c/d}{\draw[pcflow] (\p.west) -- (\q.east);}
  \node[gp] (g1) at (-1.8,-1.5) {\tn{לא נבדק: חילוץ ידע מראיונות המעוגן בגרסה ובהרשאה}};
  \node[gp] (g2) at (-5.4,-1.5) {\tn{לא נבדק: הפקה מעוגנת עם אישור מומחה ומדידת זמן מול בסיס ידני}};
  \node[gp] (g3) at (-9.0,-1.5) {\tn{לא נבדק: תרחיש \latin{3D} שנוצר מתצורת \latin{JSON} ומתעדכן עם שינוי גרסה}};
  \foreach \x in {-1.8,-5.4,-9.0}{\draw[pcBad,dashed,thick] (\x,0) -- (\x,-1.5);}
\end{pctikz}}}
\caption[מפת הספרות: ארבע תמות והחיבורים שלא נבדקו]{מפת הספרות: מספר המאמרים בכל תמה (לפי שיוך עיקרי; סך הכול \latin{20}) והחיבור שלא נבדק בין תמות עוקבות. מקור: סינתזת הסקירה בטבלה~\ref{tab:themes} ובמטריצה בטבלה~\ref{tab:litmatrix}}\label{fig:litmap}
\end{figure}
הסקירה גם מספקת שיטות הערכה שאפשר לאמץ ולא רק ממצאים. טבלה~\ref{tab:adopted} מפרטת כל שיטה, מה היא מודדת לפי מקורה, והיכן היא משמשת בפרויקט; לפי
הגישה של פרק~\ref{ch:problem}, אף שיטה אינה מוחלפת בשיפוט אנושי של שני מעריכים, והיא משמשת כבדיקה משלימה בלבד.
\begin{table}[!htbp]
\centering
\caption{שיטות הערכה שאומצו מהספרות: מה נמדד לפי המקור והשימוש בפרויקט}\label{tab:adopted}
\begin{pctable}{H{0.19},C{0.10},H{0.25},H{0.28},C{0.12}}
  \pchead{שיטה או מדד & מקור & מה היא מודדת (לפי המקור) & שימוש בפרויקט & דרישה}
  \pcrow{נאמנות לראיות ורלוונטיות ללא תשובת אמת & \cite{es2024ragas} & מדדים לכל שאילתה בלי תשובת אמת, ובהם נאמנות לראיות ורלוונטיות & \mth{G=N_{supported}/N_{claims}}; שיפוט אנושי של שני מעריכים, שופט אוטומטי רק משלים & \req{03}, \hyp{3}}
  \pcrow{שופט אוטומטי מכויל & \cite{saad2024ares} & בדיקת רכיבי \latin{RAG} בנפרד בשופט המכויל בשיפוט אנושי & בדיקה נפרדת של אחזור והפקה, בכיול על מדגם מתויג & \idlist{REQ-03, 05}}
  \pcrow{מבחן שנוצר אוטומטית ותורת התגובה לפריט & \cite{guinet2024exam} & מבחן רב-ברירה מקורפוס המשימה, וזיהוי שאלות בעלות מידע נמוך & מחוון ידני בן ארבעה ממדים; \latin{IRT} לא יופעל במדגם הקטן & \req{12}, \hyp{3}}
  \pcrow{הערכה יזומה וקטגוריות ליקוי & \cite{hua2025civask} & שלוש קבוצות ליקויים בעוזר הלמידה & \mth{D_c} ו-\mth{HR@k} בניסוח שלנו (סעיף~\ref{sec:flagship}) & \idlist{REQ-02, 05}}
  \pcrow{חוסן לשינוי ידע & \cite{zheng2025knowshift} & התנהגות \latin{RAG} כשידע הספר שונה מידע המודל & חמישה שינויי גרסה מדומים \latin{(D5)} וחסימת פריטים תלויים & \req{11}}
  \pcrow{דליפה ממאגר פרטי & \cite{zeng2024good} & דליפה באמצעות שאילתות מתוכננות & בדיקת דליפה בכל ריצה ובקרת הרשאה & \idlist{REQ-09, 10}}
  \pcrow{גודל אפקט בהכשרה & \cite{bodding2025mr} & גדלי אפקט \mth{d} למציאות מעורבת בהכשרה מקצועית & הפרש שיעור השלמה בנקודות אחוז מול \twoD{}; גישושי, ללא טענת עוצמה סטטיסטית & \hyp{4}}
\end{pctable}
\end{table}

\section{פרוטוקול החיפוש}\label{sec:protocol}
הסקירה מבוססת על חיפוש ממוקד בדפי פרסום ובמטא-דאטה רשמיים של מאמרים משנים \rng{2023}{2026} (ומקור יסוד אחד מ-\latin{2020}).
זו אינה סקירה שיטתית לפי \latin{PRISMA}: החיפוש היה ממוקד-נושא, מספר התוצאות הגולמיות לא נספר, ולכן הוא מסומן כך בטבלה
ואינו מוצג כאילו נמדד. עמודת “אחרי קריאה מלאה” מציגה את המקורות שנקראו במלואם (טבלה~\ref{tab:search}); מהם נבחרו \latin{20}
לסינתזה בסעיף~\ref{sec:review}, והיתר מוזכרים כרקע במבוא ובפרק התכן.
\begin{table}[!htbp]
\centering
\caption{פרוטוקול החיפוש: מקורות שנקראו במלואם לפי מאגר}\label{tab:search}
\begin{pctable}{H{0.17},H{0.31},C{0.10},C{0.13},C{0.12},C{0.17}}
  \pchead{מאגר & מחרוזת חיפוש & תאריך & תוצאות & אחרי כותרת ותקציר & אחרי קריאה מלאה}
  \pcrow{\latin{ACL Anthology} & \latin{RAG evaluation; long documents; tables; privacy} & \latin{11-2026} & לא נספר (חיפוש ממוקד) & --- & \latin{18}}
  \pcrow{מאגרי מו"לים \latin{(Wiley, Elsevier, Springer, IEEE, ACM)} & \latin{tacit knowledge AND LLM; mixed reality AND vocational training; XR AND aircraft} & \latin{11-2026} & לא נספר (חיפוש ממוקד) & --- & \latin{10}}
  \pcrow{כנסים נוספים \latin{(PMLR, NeurIPS, PHM, HICSS)} & \latin{exam generation RAG; technical manuals} & \latin{11-2026} & לא נספר (חיפוש ממוקד) & --- & \latin{4}}
  \pcrow{\latin{Google Scholar} & כדור שלג ואימות זהות פרסום & \latin{11-2026} & אימות פרטני & --- & \latin{0} (אימות בלבד)}
  \midrule
  \pcrow{\textbf{סה"כ (\rng{2023}{2026})} & & & & & \textbf{\latin{32}}}
\end{pctable}
\end{table}
\textbf{תנאי הכללה והחרגה:} מאמר שפיט; שנים \rng{2023}{2026} (מקור יסוד ישן יותר רק כרקע); ניסוי, תכן או סקירה שיטתית;
שפה אנגלית; הלימה לאחת מארבע התמות. \textbf{כדור שלג:} מעקב אחורה וקדימה ממאמר הדגל וממסגרות ההערכה \latin{RAGAs} ו-\latin{ARES}.
מקורות היצרנים, התקנים ודוחות השוק אינם נכללים בספירה.

\subsection*{איומים על תוקף הסקירה}
המסקנות על הפער תקפות רק ביחס לסט שנסקר, ולכן יש לקרוא אותן בזהירות:
\begin{itemize}
  \item \textbf{היעדר חיפוש שיטתי מלא:} החיפוש היה ממוקד-נושא ובמספר מאגרים מוגבל, ולכן ייתכנו מחקרים רלוונטיים שלא אותרו. “לא נמצא” אינו “לא קיים”.
  \item \textbf{הטיית בחירה:} עשרים המאמרים נבחרו לפי הלימה לארבע התמות; מאמרים שאינם מתאימים לשרשרת הערך עלולים להיעדר, וכך להעצים את מראה הפער.
  \item \textbf{שפה ושדה מתפתח:} נכללו רק מאמרים באנגלית, ובתחום משתנה במהירות: מחקרים שיפורסמו אחרי נובמבר \latin{2026} עשויים לסגור חלק מהפער.
  \item \textbf{העברה בין תחומים:} חלק מהמחקרים נערכו בהוראה אקדמית ובייצור ולא בתעשיית התעופה והחלל; תקפות ההעברה אינה מובטחת.
  \item \textbf{מקורות אפורים:} דפי יצרנים ודוחות שוק לא נכללו בסינתזה, והם נבחנים בנפרד בפרק~\ref{ch:market} עם סימון המעמד שלהם.
\end{itemize}

\section{מטריצת ספרות}\label{sec:matrix}
המטריצה בטבלה~\ref{tab:litmatrix} בודקת שישה ממדים: הערכת \latin{RAG}, מסמכים וגרסאות, ידע סמוי, \threeD{} ומציאות מורחבת,
למידה ושאלות, ואבטחה ופרטיות. נקודה מסמנת התייחסות משמעותית; תא ריק משמעו שהממד אינו מוקד המחקר, ולא שהחוקרים הוכיחו
היעדר יכולת.
\begin{table}[!htbp]
\centering
\caption{מטריצת מושגים בשישה ממדים (\latin{20} מאמרים): אף מחקר אינו מכסה את כל השורה}\label{tab:litmatrix}
\begin{pctable}{c,C{0.10},C{0.13},C{0.13},C{0.13},C{0.13},C{0.13},C{0.13}}
  \pchead{מאמר & שנה & הערכת \latin{RAG} & מסמכים וגרסאות & ידע סמוי & \threeD{} \latin{/XR} & למידה ושאלות & אבטחה}
  \pcrow{\cite{igoa2024transfer} & \latin{2024} & & & \pcdot & & & }
  \pcrow{\cite{finkel2026unlocking} & \latin{2026} & & & \pcdot & & & }
  \pcrow{\cite{zuin2025tacit} & \latin{2025} & & & \pcdot & & & }
  \pcrow{\cite{ji2023survey} & \latin{2023} & \pcdot & & & & & }
  \pcrow{\cite{barnett2024seven} & \latin{2024} & \pcdot & \pcdot & & & & }
  \pcrow{\cite{es2024ragas} & \latin{2024} & \pcdot & & & & & }
  \pcrow{\cite{saad2024ares} & \latin{2024} & \pcdot & & & & & }
  \pcrow{\cite{tao2025treerag} & \latin{2025} & \pcdot & \pcdot & & & & }
  \pcrow{\cite{zheng2025knowshift} & \latin{2025} & \pcdot & \pcdot & & & \pcdot & }
  \pcrow{\cite{lukens2025manuals} & \latin{2025} & \pcdot & \pcdot & & & & }
  \pcrow{\cite{zeng2024good} & \latin{2024} & \pcdot & & & & & \pcdot}
  \pcrow{\cite{hua2025civask} & \latin{2025} & \pcdot & \pcdot & & & \pcdot & }
  \pcrow{\cite{guinet2024exam} & \latin{2024} & \pcdot & & & & \pcdot & }
  \pcrow{\cite{ramkumar2026pedagogical} & \latin{2026} & & \pcdot & & & \pcdot & }
  \pcrow{\cite{pradeesh2025mcq} & \latin{2025} & & \pcdot & & & \pcdot & }
  \pcrow{\cite{lee2025impact} & \latin{2025} & & & & & \pcdot & }
  \pcrow{\cite{bodding2025mr} & \latin{2025} & & & & \pcdot & \pcdot & }
  \pcrow{\cite{somerville2025xr} & \latin{2025} & & & & \pcdot & \pcdot & }
  \pcrow{\cite{jalilvand2024xr} & \latin{2024} & & & & \pcdot & \pcdot & }
  \pcrow{\cite{yang2023xr} & \latin{2023} & & & & \pcdot & & }
  \midrule
  \pcrow{\textbf{הפרויקט שלנו} & \latin{2027} & \pcdot & \pcdot & \pcdot & \pcdot & \pcdot & \pcdot}
\end{pctable}
\end{table}
המטריצה מראה שלוש קבוצות כמעט נפרדות: מחקרי ידע סמוי (תמה 1) אינם בודקים אחזור או תרגול; מחקרי \latin{RAG} (תמה 2) מתמקדים
בהערכת אחזור ומסמכים, ורק מחקר אחד בוחן אבטחה; ומחקרי ה\threeD{} (תמה 4) אינם מעגנים את תוכן ההדרכה במסמכים. המחקרים המכסים
את מירב הממדים \cite{zheng2025knowshift,hua2025civask} אינם כוללים ידע סמוי, תרגול \latin{3D}, אישור מומחה או מדידת זמן מול בסיס
ידני --- השילוב שהפרויקט מבקש לבחון.

\section{בחירת מאמר הדגל}\label{sec:flagshipchoice}
מאמר הדגל שנבחר הוא \latin{Hua} ועמיתיו \latin{(2025)}, על הערכה יזומה של עוזר למידה מבוסס \latin{RAG} בהוראה הנדסית
\cite{hua2025civask}. הוא עדיף לפרויקט על המטא-אנליזה של מציאות מעורבת \cite{bodding2025mr}, משום שהוא מציג מערכת שאפשר
לשחזר את שלביה: נתונים, אחזור, שאלות וניסוי (טבלה~\ref{tab:flagship}). המטא-אנליזה משמשת עוגן להשערת ה\threeD{}, ומאמרם של
\latin{Lukens} ועמיתיו משמש בקרה משלימה בהקשר של מדריכי תחזוקה \cite{lukens2025manuals}.
\begin{table}[!htbp]
\centering
\caption{בחירה מנומקת של מאמר הדגל}\label{tab:flagship}
\begin{pctable}{H{0.24},C{0.18},C{0.18},H{0.40}}
  \pchead{קריטריון & \latin{CivASK (2025)} & מטא-אנליזת \latin{MR (2025)} & נימוק}
  \pcrow{במה מחקרית ועדכניות & כתב עת שפיט, \latin{2025} & כתב עת שפיט, \latin{2025} & שניהם עומדים בתנאי הסף}
  \pcrow{קרבה לבעיה & אחזור ידע, למידה והערכה & אפקטיביות הדרכה במציאות מעורבת & \latin{CivASK} עוסק בליבת הפתרון: ידע מאוחזר שהופך ללמידה}
  \pcrow{חזרתיות & תהליך ומדגם ניסוי מתוארים & סינתזה של \latin{53} מחקרים & אפשר להעתיק את מסגרת ההערכה של \latin{CivASK} לקורפוס תחזוקה}
  \pcrow{מדדים כמותיים & \num{742} שאילתות, \num{374} שאלות & גדלי אפקט \mth{d} & \latin{CivASK} מספק מערך מדידה שאפשר לאמץ כמסלול גיבוי}
\end{pctable}
\end{table}

\section{ניתוח מאמר הדגל}\label{sec:flagship}
\subsection*{מטרה ושאלה}
המאמר מבקש לזהות ליקויים בתמיכה של עוזר למידה הנדסי המבוסס על אחזור מסמכי קורס, בלי להסתמך רק על דיווחי משתמשים
\cite{hua2025civask}. כלומר: הערכה \textbf{יזומה} של המערכת לפני שכשליה מתגלים בשימוש.
\subsection*{שיטה, נתונים וסביבת הניסוי}
במאמר נותחו \num{742} שאילתות של משתמשים, ונוצרו מהן \num{374} שאלות לבדיקת העוזר. אלה מספרי המדגם של המאמר ולא של הפיילוט
שלנו. הקורפוס הוא חומרי הוראה בהנדסה, והמערכת נבדקת מול שאלות שנגזרו משאילתות אמיתיות, כך שהסביבה מייצגת שימוש ולא רק
אוסף שאלות מלאכותי. הארכיטקטורה הנבדקת: ייצוג חומרי ההוראה במאגר ידע, אחזור קטעים רלוונטיים לשאלת הלומד, יצירת תשובה
בידי מודל שפה, והערכת התשובה מול ידע המקור.
\subsection*{ניסוח מתמטי ומדדים}
הרכיב המרכזי הוא מודל \latin{RAG} הסטנדרטי שעליו נשענת הארכיטקטורה \cite{lewis2020retrieval}: לשאלה \mth{x} מאוחזרים
\mth{k} המקטעים \mth{z} בעלי הציון הגבוה במאגר, והתשובה \mth{y} נוצרת בהתניה עליהם,
\begin{equation}
  p(y\mid x)\;\approx\;\sum_{z\,\in\,\mathrm{top}\text{-}k\left(p_\eta(\cdot\mid x)\right)} p_\eta(z\mid x)\;p_\theta(y\mid x,z),
\end{equation}
כאשר \mth{p_\eta} הוא המאחזר ו-\mth{p_\theta} הוא מודל השפה. מערך ההערכה של המאמר מסווג ליקויים לקטגוריות. את המדדים
נבטא בפרויקט באופן פורמלי כך, וזוהי \textbf{הפורמליזציה שלנו} לקטגוריות המדווחות ולא ציטוט משוואות מתוך המאמר:
\begin{equation}
  D_c=\frac{\bigl|\{\,q\in Q:\ \text{\texthebrew{ליקוי מסוג}}\ c\,\}\bigr|}{|Q|},\qquad
  HR@k=\frac{1}{|Q|}\sum_{q\in Q}\mathbf{1}\bigl[z^{*}_{q}\in \mathrm{top}\text{-}k(q)\bigr],
\end{equation}
כאשר \mth{Q} הוא מערך שאלות הבדיקה, \mth{c} קטגוריית ליקוי (פער במאגר, אי-הבנת הקשר, אחזור חסר), ו-\mth{z^{*}_{q}} הקטע הנכון לשאלה
\mth{q}. \mth{D_c} הוא שיעור הליקויים בקטגוריה, ו-\mth{HR@k} הוא שיעור השאלות שהקטע הנכון נמצא בראש \mth{k} התוצאות. כך אפשר להשוות
בין סבבי בדיקה ולזהות איזו קטגוריה ירדה לאחר תיקון. בפיילוט נוסיף את הנאמנות \mth{G=N_{supported}/N_{claims}} ואת חיסכון הזמן
\mth{1-T_{RAG}/T_{manual}} (סעיף~\ref{sec:soleval}), שאינם חלק ממסגרת המאמר.
\subsection*{ממצאים}
המסגרת איתרה שלוש קבוצות ליקויים בעוזר הלמידה: פערים במאגר הידע, אי-הבנה של הקשר השאלה ואחזור חסר
\cite{hua2025civask}. כל אחת מהן מקבילה לסיכון בפרויקט: פער במאגר (\risk{01}), תוכן שאינו נתמך במקור (\risk{03})
ואיכות אחזור נמוכה (\risk{05}).
\subsection*{מגבלות}
ככל שעולה מתיאור המאמר והיקפו, אלו המגבלות הרלוונטיות לפרויקט:
\begin{itemize}
  \item \textbf{העברה בין תחומים:} הקורפוס הוא חומר הוראה אקדמי, ואין הוכחה שאותה איכות תישמר במסמכי נוהל ותחזוקה של מטוסים.
  \item \textbf{ידע סטטי:} לא מדווחת בדיקה של התנהגות המערכת כשמסמך מתעדכן, כלומר שליטה בגרסאות.
  \item \textbf{ללא שער מומחה:} לא מדווח אישור מומחה כתנאי לפרסום, ולכן אין בסיס להצהיר על אמינות מקצועית.
  \item \textbf{מדדי יעילות:} לא מדווח חיסכון בזמן ההכנה לעומת תהליך ידני, ולא נבדק ידע סמוי מראיונות.
  \item \textbf{ללא תרגול מרחבי ופריסה מבודדת:} אין שכבת סימולציה, ולא מדווח ניסוי תחת אילוץ של עיבוד מקומי בלבד.
\end{itemize}
\subsection*{הקשר לפרויקט}
נאמץ את מבנה ההערכה (שאלות אמת, בדיקת אחזור, אישור שאלות ומדד \mth{D_c}) ונרחיב אותו בארבעה מקומות: עיגון כל טענה בגרסה
ובעמוד, אישור מומחה לפני פרסום, מדידת זמן ההפקה מול בסיס ידני, ותרגול \latin{3D} הנבדק מול \latin{2D}. אין להעתיק
מן המאמר תוצאות שלא נמדדו בו.
טבלה~\ref{tab:adoptext} מפרידה בין מה שמאומץ מן המאמר כפי שהוא לבין מה שהפרויקט מוסיף; ההשוואה נשענת על תיאור המאמר והיקפו, וכל תא של “לא מדווח” פירושו
שהמידע אינו מופיע בתיאור שנסקר, ולא שהמאמר מוכיח היעדר.
\begin{table}[!htbp]
\centering
\caption{מאמר הדגל מול הפרויקט: מה מאומץ ומה מורחב}\label{tab:adoptext}
\begin{pctable}{H{0.21},H{0.27},H{0.34},C{0.12}}
  \pchead{רכיב & במאמר הדגל \cite{hua2025civask} & בפרויקט & סוג}
  \pcrow{הערכה יזומה לפני שהכשלים מתגלים בשימוש & מטרה מרכזית של המאמר & בדיקת המערכת לפני פרסום כל יחידה & אימוץ}
  \pcrow{שאלות בדיקה & \num{374} שאלות שנוצרו מתוך \num{742} שאילתות משתמשים & לפחות \latin{20} שאלות אמת ו-\latin{10} מקרי חוסר-ראיות שמכין מעריך, כי בפיילוט אין משתמשים & התאמה}
  \pcrow{קטגוריות ליקוי ושיעור אחזור & שלוש קבוצות ליקויים & אותן קטגוריות כ-\mth{D_c} ו-\mth{HR@k}, ובנוסף נאמנות \mth{G} & אימוץ}
  \pcrow{גרסאות מסמך & לא מדווח & חמישה שינויי גרסה וחסימת פריטים תלויים (\req{11}) & הרחבה}
  \pcrow{אישור מומחה & לא מדווח כתנאי לפרסום & שער חובה לפני פרסום (\req{06}) & הרחבה}
  \pcrow{זמן הפקה & לא מדווח & חציון זמן מול בסיס ידני (\req{04}) & הרחבה}
  \pcrow{תרגול \threeD{} ופריסה מקומית & לא מדווח & תבנית \threeD{} מתצורת \latin{JSON}; אפס בקשות לרשת חוץ (\req{09}) & הרחבה}
\end{pctable}
\end{table}

\section{פער הידע}\label{sec:gap}
הפער נגזר ישירות מהמטריצה (טבלה~\ref{tab:litmatrix}) ומהסינתזה (טבלה~\ref{tab:themes}). הוא מתורגם לניסוי בחמישה רכיבים:
(א) דיוק חילוץ ידע מראיונות (\req{16}, \hyp{1}); (ב) זמן יצירה ואישור של יחידה (\hyp{2}); (ג) איכות אחזור ותמיכה בעדויות
ואיכות שאלות (\hyp{3}); (ד) תרגול \latin{3D} מול \latin{2D} (\hyp{4}); (ה) עמידה באילוץ הפריסה המקומית ובבדיקת דליפה
(\req{09}). בכך ההשערות בפרק~\ref{ch:problem} נובעות ישירות מן הספרות.
\begin{pcguide}[בדיקה]
המשפט עומד בשלושה תנאים: נשען על מקור, מתאים לתאים הריקים במטריצה, ומוביל ישירות לשאלת המחקר
ולחדשנות (פרק~\ref{ch:alt}).
\end{pcguide}
\begin{pcbox}[pcBad]{פער הידע (\latin{Engineering Gap})}
עד כה לא נבדק, בסביבה מקומית מאובטחת, מנגנון אחד המחבר בין חילוץ ידע סמוי מראיונות מומחים, הפקת יחידות הדרכה
ומבחנים מעוגנים ומגורסאים ממסמכים הנדסיים, ותצורות תרחיש \threeDadj{} שנוצרות מהן אוטומטית ומתעדכנות עם שינוי גרסה, תוך מדידת
דיוק החילוץ, זמן ההפקה והנאמנות העובדתית מול בסיס ידני ומול תרגיל \mbox{דו-ממדי} זהה.
\end{pcbox}
\pcChapterEnd{lit}
